%=============================================================================!
% 4.9  COUPLED OSCILLATORS AND NORMAL MODES                                   !
%=============================================================================!
\section{Coupled oscillators and normal modes}
\label{sec:os-modes}

For small displacements from equilibrium, the forces in a molecule or
crystal can be approximated by coupled springs. When one atom moves, its
neighbours feel it, so the motion of one atom cannot be solved on its own.
Within this harmonic approximation the system has simple collective
motions, its normal modes, in which the atoms oscillate at a common
angular frequency. Their sum describes a general motion of the coupled
system. Finding the modes needs two new tools, matrices and their
eigenvalues.

\subsection*{Two coupled carts}

Two carts of mass $m$ stand on a smooth track between two walls, each
joined to its wall by a spring of stiffness $k$, and to each other by a
third such spring; all three springs are at their natural lengths when
the carts rest. Let $x_1$ and $x_2$ be the carts' displacements from rest.
The left spring is stretched by $x_1$ and pulls cart 1 with $-kx_1$; the
middle spring is stretched by $x_2 - x_1$ and pulls cart 1 forwards with
$k(x_2 - x_1)$ and cart 2 backwards with $-k(x_2 - x_1)$; the right spring
is squashed by $x_2$ and pushes cart 2 with $-kx_2$. Adding the forces on
each cart,
\begin{equation}
   m\ddot{x}_1 = -2k\,x_1 + k\,x_2 , \qquad
   m\ddot{x}_2 = k\,x_1 - 2k\,x_2 .
   \label{eq:os-two-carts}
\end{equation}
Each equation contains both displacements: the carts are coupled. Started
by pulling one cart aside, the swing passes back and forth between them
(\cref{fig:os-modes}c), and neither moves as a single spring. A table of
numbers organises such equations.

\begin{toolbox}[Matrices]
A matrix is a rectangular table of numbers. A square matrix $\vb{A}$
with entries $A_{ab}$, in row $a$ and column $b$, acts on a vector
$\vb{u} = (u_1, \dots, u_n)$ to give the vector $\vb{A}\vb{u}$ whose
$a$th component is the scalar product of row $a$ with $\vb{u}$,
\begin{equation*}
   (\vb{A}\vb{u})_a = \sum_{b=1}^{n} A_{ab}\,u_b ;
   \qquad\text{for example}\quad
   \begin{pmatrix} 2 & -1\\ -1 & 2 \end{pmatrix}
   \begin{pmatrix} u_1\\ u_2 \end{pmatrix}
   = \begin{pmatrix} 2u_1 - u_2\\ -u_1 + 2u_2 \end{pmatrix} .
\end{equation*}
This product is linear: $\vb{A}(\vb{u} + \vb{w}) = \vb{A}\vb{u} +
\vb{A}\vb{w}$ and $\vb{A}(c\vb{u}) = c\,\vb{A}\vb{u}$ for a
number $c$, since each component is a sum of products. A number times a
matrix multiplies every entry, so $(c\vb{A})\vb{u} = c\,(\vb{A}\vb{u})$.
A diagonal matrix $\vb{B}$ has zeros everywhere off its diagonal, the
entries $B_{aa}$ from top left to bottom right, so it multiplies each
component by its own entry, $(\vb{B}\vb{u})_a = B_{aa}u_a$. A matrix with
$A_{ab} = A_{ba}$ for every pair is symmetric, and the identity matrix
$\vb{I}$, with ones on the diagonal and zeros elsewhere, leaves every
vector unchanged, $\vb{I}\vb{u} = \vb{u}$.
\end{toolbox}

With the displacements as a vector $\vb{x} = (x_1, x_2)$,
\cref{eq:os-two-carts} reads
\begin{equation*}
   m\ddot{\vb{x}} = -\boldsymbol{\Phi}\vb{x} , \qquad
   \boldsymbol{\Phi} = k\begin{pmatrix} 2 & -1\\ -1 & 2 \end{pmatrix} .
\end{equation*}
The matrix $\boldsymbol{\Phi}$, called the force-constant matrix, is the
table of second derivatives of the energy stored in the springs, $U =
\frac12 kx_1^2 + \frac12 k(x_2 - x_1)^2 + \frac12 kx_2^2$. By the chain
rule, $\partial U/\partial x_1 = kx_1 - k(x_2 - x_1) = 2kx_1 - kx_2$,
minus the force on cart 1, so $\partial^2U/\partial x_1^2 = 2k$ and
$\partial^2U/\partial x_2\partial x_1 = -k$, the first row of
$\boldsymbol{\Phi}$; the second row follows in the same way. Such a table
of all the second derivatives of an energy is called its Hessian, the
table promised in \cref{sec:en-curl}, and the equality of mixed partial
derivatives proved there makes it symmetric.

\subsection*{Normal modes}

Look for a motion in which both carts swing as sinusoids with the same
angular frequency and phase, $\vb{x} = \vb{u}\cos(\omega t + \phi)$, with
$\vb{u}$ a fixed vector, the pattern of the motion. Then $\ddot{\vb{x}} =
-\omega^2\vb{x}$, and the equation of motion becomes $-m\omega^2\vb{u}
\cos(\omega t + \phi) = -\boldsymbol{\Phi}\vb{u}\cos(\omega t + \phi)$, which holds
at all times only if
\begin{equation}
   \boldsymbol{\Phi}\vb{u} = m\omega^2\,\vb{u} .
   \label{eq:os-eigen}
\end{equation}
The matrix must turn the pattern into a multiple of itself. A motion of
this kind is called a normal mode, and \cref{eq:os-eigen} is an eigenvalue
problem.

\begin{toolbox}[Eigenvalues and eigenvectors]
A non-zero vector $\vb{u}$ with $\vb{A}\vb{u} = \lambda\vb{u}$ for some
number $\lambda$ is an eigenvector of $\vb{A}$, and $\lambda$ its
eigenvalue. For a $2\times2$ matrix with entries $a$, $b$ in the first
row and $c$, $d$ in the second, the condition is the pair of equations
\begin{equation*}
   (a - \lambda)\,u_1 + b\,u_2 = 0 , \qquad
   c\,u_1 + (d - \lambda)\,u_2 = 0 .
\end{equation*}
Multiply the first by $d - \lambda$ and the second by $b$, and subtract:
the $u_2$ terms cancel, leaving $[(a - \lambda)(d - \lambda) - bc]\,u_1 =
0$. Multiplying instead by $c$ and by $a - \lambda$ and subtracting leaves
the same bracket times $u_2$. Since $u_1$ and $u_2$ are not both zero,
the bracket must vanish:
\begin{equation*}
   (a - \lambda)(d - \lambda) - bc = 0 .
\end{equation*}
This quadratic equation, the characteristic equation, gives the
eigenvalues, and putting each back into either of the pair gives its
eigenvector, up to a factor: any multiple $c\vb{u}$ of an eigenvector is
another, since $\vb{A}(c\vb{u}) = c\,\vb{A}\vb{u} = \lambda\,(c\vb{u})$.
For a symmetric matrix, $c = b$, multiplying out gives $\lambda^2 - (a +
d)\lambda + (ad - b^2) = 0$, whose discriminant, the quantity under the
square root in the formula for the roots of a quadratic, is
\begin{equation*}
   (a + d)^2 - 4(ad - b^2) = a^2 - 2ad + d^2 + 4b^2 = (a - d)^2 + 4b^2 ,
\end{equation*}
expanding $(a + d)^2$ and collecting $2ad - 4ad = -2ad$. A sum of squares
is never negative, so the eigenvalues are real. An $n\times
n$ symmetric matrix likewise has $n$ real eigenvalues, with eigenvectors
that can be chosen at right angles to one another; we use this general
result without proof.
\end{toolbox}

For the carts, dividing \cref{eq:os-eigen} by $k$ gives the eigenvalue
problem, with $\lambda = m\omega^2/k$, of the matrix with rows $(2, -1)$ and $(-1, 2)$, whose
characteristic equation is $(2 - \lambda)^2 - 1 = 0$, so $2 - \lambda =
\pm1$ and $\lambda = 1$ or 3.
\begin{itemize}
\item $\lambda = 1$: the first equation, $(2 - 1)u_1 - u_2 = 0$, gives
   $\vb{u} = (1, 1)$. Both carts move together, the middle spring is never
   stretched, and $\omega = \sqrt{k/m}$.
\item $\lambda = 3$: $(2 - 3)u_1 - u_2 = 0$ gives $\vb{u} = (1, -1)$. The
   carts move in opposite directions, and $\omega = \sqrt{3k/m}$.
\end{itemize}
With $m = \SI{1}{\kilogram}$ and $k = \SI{100}{\newton\per\metre}$ the
two modes have $\omega = 10$ and $\SI{17.32}{\radian\per\second}$
(\cref{fig:os-modes}a, b). The two patterns are at right angles,
$(1, 1)\cdot(1, -1) = 0$.

Because the equations are linear, a sum of normal modes is also a motion,
as in \cref{sec:os-driven}. Any starting displacements split between the
two patterns, $(x_1, x_2) = \frac12(x_1 + x_2)\,(1, 1) + \frac12(x_1 -
x_2)\,(1, -1)$, and the starting velocities likewise, so an amplitude and
a phase for each mode can match any starting state, and by the uniqueness
of solutions (\cref{sec:ne-equations}) every motion of the carts is a sum
of normal modes. (A mode with $\omega = 0$, met below, contributes a
steady drift $(c_1 + c_2t)\,\vb{u}$ in place of a sinusoid.) Pulling cart
1 out by \SI{4}{\centi\metre} and releasing both carts at rest starts the
two modes equally, since $(4, 0) = 2\,(1, 1) + 2\,(1, -1)$, each at a
turning point, with phase zero:
\begin{equation*}
   \vb{x}(t) = 2\,(1, 1)\cos 10t + 2\,(1, -1)\cos 17.32t
   \quad\text{in centimetres.}
\end{equation*}
The fast mode gains on the slow one by $17.32 - 10 = 7.32$ radians every
second. When it has gained $\pi$, its cosine is minus the slow one's, so
cart 1 stands nearly still and cart 2 swings \SI{4}{\centi\metre} out;
when it has gained $2\pi$ the two are back in step. The swing therefore
passes from one cart to the other and back every $2\pi/7.32 =
\SI{0.858}{\second}$ (\cref{fig:os-modes}c).

\begin{figure}[htb]
\centering
\includegraphics{ch04_oscillations/modes}
\caption{Two carts of \SI{1}{\kilogram} between walls, joined by three
springs of \SI{100}{\newton\per\metre}. (a) The slow mode: both start
\SI{2}{\centi\metre} out and swing together at
\SI{10}{\radian\per\second}. (b) The fast mode: cart
1 starts \SI{2}{\centi\metre} out one way and cart 2 \SI{2}{\centi\metre}
the other, and they swing against each other at
\SI{17.32}{\radian\per\second}. (c) Only cart 1 displaced: the sum of the
two modes, in which the swing passes back and forth between the carts.
Cart 1 is solid and cart 2 dashed, so both remain identifiable where
they move together.}
\label{fig:os-modes}
\end{figure}

\subsection*{Many atoms}

Near a minimum of the potential energy $U(\vb{r}^N)$, let $q_1, \dots,
q_n$, with $n = 3N$, be the displacements of the coordinates from the
minimum. Along the straight line from the minimum in the direction of a
displacement $\vb{q}$, the energy is $f(s) = U(\text{minimum} + s\vb{q})$.
The chain rule along a path of \cref{sec:en-gradient}, with $\dd(\text{minimum} +
s\vb{q})/\dd s = \vb{q}$, gives $f'(s) = \sum_a(\partial U/\partial
q_a)\,q_a$. Each $\partial U/\partial q_a$ is itself a function of
position, so the same rule gives its rate of change along the line as
$\sum_b(\partial^2U/\partial q_b\,\partial q_a)\,q_b$, and since the $q_a$
are fixed numbers,
\begin{equation*}
   f''(s) = \sum_a\sum_b\frac{\partial^2U}{\partial q_b\,\partial q_a}\,
   q_aq_b = \sum_{a,b}\Phi_{ab}\,q_aq_b ,
\end{equation*}
where $\sum_{a,b}$ runs over every pair $a$, $b$, $n^2$ terms in all, and
$\Phi_{ab} = \partial^2U/\partial q_a\partial q_b$, the same in either
order (\cref{sec:en-curl}), is the Hessian. At the minimum $f'(0) = 0$, so
the Taylor series of $f$ about $s = 0$, stopped after the term in $s^2$
and taken at $s = 1$, gives
\begin{equation*}
   U = U_{\min} + \tfrac12\sum_{a,b} \Phi_{ab}\,q_a q_b + \order{q^3} ,
\end{equation*}
with $\Phi_{ab}$ taken at the minimum and $q$ the length of $\vb{q}$. The
step $s = 1$ is not short, but by \cref{sec:mo-taylor} the error is at
most the largest size of $f'''$ divided by $3! = 6$, and $f'''$, from the
chain rule a third time, is a sum of third derivatives of $U$ times
products of three components of $\vb{q}$, so the error shrinks as $q^3$.
This is the many-dimensional version of \cref{sec:os-small}.

The same expansion gives the test promised in \cref{sec:en-gradient}. Let
$\vb{u}_1, \dots, \vb{u}_n$ be eigenvectors of $\boldsymbol{\Phi}$ of
length 1, at right angles to one another, with eigenvalues $\lambda_1,
\dots, \lambda_n$, and write a displacement as $\vb{q} = \sum_k
c_k\vb{u}_k$. Since $\sum_{a,b}\Phi_{ab}q_aq_b = \vb{q}\cdot
(\boldsymbol{\Phi}\vb{q})$, and $\vb{u}_k\cdot\vb{u}_l$ is 1 for $k = l$
and 0 otherwise, the quadratic term is $\frac12\sum_k\lambda_k c_k^2$.
Where the gradient vanishes, $U$ therefore has a minimum if every
eigenvalue of $\boldsymbol{\Phi}$ is positive, a maximum if every one is
negative, and a saddle point if there are some of each. Zero eigenvalues
need separate attention: higher terms may decide the behaviour, or a
symmetry may allow free motion, as in the translation below.

Differentiating the expansion with respect to $q_c$, with the other
displacements held fixed, the product rule gives each term $\Phi_{ab}q_aq_b$
the derivative $\Phi_{cb}q_b$ if $a = c$, plus $\Phi_{ac}q_a$ if $b = c$,
and nothing if neither index is $c$, so
\begin{equation*}
   \frac{\partial U}{\partial q_c}
   = \tfrac12\sum_b \Phi_{cb}\,q_b + \tfrac12\sum_a \Phi_{ac}\,q_a
   = \sum_b \Phi_{cb}\,q_b ,
\end{equation*}
the second step renaming $a$ as $b$ and using $\Phi_{bc} = \Phi_{cb}$.
Neglecting terms of order $q^2$, the force is therefore $F_c = -\sum_b
\Phi_{cb}q_b$, and Newton's law for each coordinate, $m_c\ddot{q}_c = F_c$,
reads $\vb{M}\ddot{\vb{q}} = -\boldsymbol{\Phi}\vb{q}$, with $\vb{M}$ the
diagonal matrix of the masses, each atom's mass appearing for its $x$,
$y$ and $z$. A normal mode $\vb{q} = \vb{u}\cos(\omega t + \phi)$ now needs
$\boldsymbol{\Phi}\vb{u} = \omega^2\vb{M}\vb{u}$. When the masses differ, this
is not yet an eigenvalue problem of a symmetric matrix, but the
substitution $w_a = \sqrt{m_a}\,u_a$ makes it one. The $a$th row reads
$\sum_b \Phi_{ab}u_b = \omega^2 m_a u_a$; putting $u_b = w_b/\sqrt{m_b}$
on the left and $m_au_a = \sqrt{m_a}\,w_a$ on the right, and dividing by
$\sqrt{m_a}$,
\begin{equation*}
   \sum_b \frac{\Phi_{ab}}{\sqrt{m_a m_b}}\,w_b = \omega^2 w_a ,
\end{equation*}
the eigenvalue problem of the mass-weighted Hessian, which is symmetric.
Its $n$ eigenvalues are the squared angular frequencies of the $n$ normal
modes. The function \texttt{normal\_modes} of \texttt{mdlab} solves it with
NumPy's routine for symmetric matrices, \texttt{eigh}:

\begin{code}
def normal_modes(
    hessian: ArrayLike, masses: ArrayLike, *, force_to_accel: float
) -> tuple[NDArray, NDArray]:
    hessian = np.asarray(hessian, dtype=float)
    root_m = np.sqrt(np.asarray(masses, dtype=float))
    weighted = hessian / np.outer(root_m, root_m)  # Φ_ab / √(m_a m_b)
    omega_squared, vectors = np.linalg.eigh(weighted)
    modes = vectors / root_m[:, np.newaxis]  # w_a / √m_a back to u_a
    return omega_squared * force_to_accel, modes
\end{code}

The lone \texttt{*} in the first line makes the argument after it a
keyword argument, passed by name as \texttt{force\_to\_accel=1.0}.
\texttt{np.outer(root\_m, root\_m)} is the table of products $\sqrt{m_a
m_b}$, so the division forms the mass-weighted Hessian entry by entry;
\texttt{eigh} returns its eigenvalues in increasing order and its
eigenvectors as columns, each of length 1 and of either sign.
\texttt{root\_m[:, np.newaxis]}, in which \texttt{:} takes the whole of
the one dimension that \texttt{root\_m} has, is the column of the
$\sqrt{m_a}$, as in \cref{sec:ne-atoms}, so the next division divides row
$a$ by $\sqrt{m_a}$ and turns each $\vb{w}$ back into a pattern of
displacements $\vb{u}$. The factor \texttt{force\_to\_accel} is 1 in SI
units and the Newton factor \num{9.648533e-3} of \cref{sec:ne-atoms} for a
Hessian in \si{\electronvolt\per\angstrom\squared} and masses in
\si{\amu}, giving $\omega^2$ in \si{\per\femto\second\squared}, as for
the lithium atom of \cref{sec:os-small}.

\subsection*{A linear molecule of three atoms}

A molecule like carbon dioxide, an oxygen atom on either side of a carbon
atom in a line, can be modelled along its axis by two springs of
stiffness $k$ between the three atoms, with masses $m_{\mathrm O}$,
$m_{\mathrm C}$, $m_{\mathrm O}$. Its energy is $U = \frac12 k(x_2 -
x_1)^2 + \frac12 k(x_3 - x_2)^2$, with no walls, and differentiating twice
as for the carts gives a Hessian with rows $k(1, -1, 0)$, $k(-1, 2, -1)$
and $k(0, -1, 1)$. It has three modes
(\cref{fig:os-triatomic}). The first, $\vb{u} = (1, 1, 1)$, moves the whole
molecule without stretching a spring, so $\boldsymbol{\Phi}\vb{u} =
\vb{0}$ and $\omega = 0$: no force resists this pattern, and the molecule
drifts along it at a steady velocity, a motion of the whole called a
translation. In the second, $\vb{u} = (1, 0, -1)$, the
carbon stays still while the oxygens move in and out together;
$\boldsymbol{\Phi}\vb{u} = k(1, 0, -1)$ and $\vb{M}\vb{u} = (m_{\mathrm O}, 0,
-m_{\mathrm O})$, so $\omega^2 = k/m_{\mathrm O}$. In the third, the oxygens
move one way and the carbon the other, so that the centre of mass stays
still, $m_{\mathrm O}u_1 + m_{\mathrm C}u_2 + m_{\mathrm O}u_3 = 0$ with
$u_1 = u_3 = 1$, which gives $\vb{u} = (1, -2m_{\mathrm O}/m_{\mathrm C},
1)$, and
\cref{ex:os-triatomic} shows that $\omega^2 = (k/m_{\mathrm O})(1 +
2m_{\mathrm O}/m_{\mathrm C})$. With $m_{\mathrm O} = \SI{15.999}{\amu}$ and
$m_{\mathrm C} = \SI{12.011}{\amu}$, the third mode is faster than the
second by $\sqrt{1 + 2\times15.999/12.011} = 1.914$, whatever $k$ is.

\begin{figure}[htb]
\centering
\begin{tikzpicture}[line cap=round, scale=1.0,
   disp/.style={-{Stealth[length=4pt]}, line width=0.9pt, accent}]
   \foreach \row/\name/\uo/\uc/\ut in {0/translation/0.5/0.5/0.5,
      1/symmetric stretch/0.5/0/-0.5, 2/antisymmetric stretch/0.3/-0.799/0.3} {
      \begin{scope}[yshift=-\row*1.0cm]
         \draw[black!40, line width=0.6pt] (0,0) -- (3.2,0);
         \fill[black!25, draw=black, line width=0.4pt] (0,0) circle (0.2);
         \fill[black!60, draw=black, line width=0.4pt] (1.6,0) circle (0.18);
         \fill[black!25, draw=black, line width=0.4pt] (3.2,0) circle (0.2);
         \ifdim \uo pt=0pt \else \draw[disp] (0,0.35) -- ++(\uo,0); \fi
         \ifdim \uc pt=0pt \else \draw[disp] (1.6,0.35) -- ++(\uc,0); \fi
         \ifdim \ut pt=0pt \else \draw[disp] (3.2,0.35) -- ++(\ut,0); \fi
         \node[anchor=west] at (4.2,0) {\small \name};
      \end{scope}
   }
   \node[below] at (0,-2.25) {\small O};
   \node[below] at (1.6,-2.25) {\small C};
   \node[below] at (3.2,-2.25) {\small O};
\end{tikzpicture}
\caption{The three normal modes of a linear O--C--O chain along its axis.
Arrows show the pattern $\vb{u}$ of displacements, to scale within each
mode. In the antisymmetric stretch the carbon moves $2m_{\mathrm O}/
m_{\mathrm C} = 2.66$ times as far as each oxygen, in the opposite
direction, so that the centre of mass stays still.}
\label{fig:os-triatomic}
\end{figure}

\notebook{04}{animate the normal modes of the carts and the triatomic
chain}

\begin{selfcheck}
In the slow mode of the two carts, why does the middle spring not affect
the angular frequency?
\end{selfcheck}
